\chapter{Introdução}

Antes de abordar os tópicos que são objetivos deste guia, é importante estabelecer alguns pontes que devem ser levados em consideração, durante a criação do seu trabalho:

\begin{enumerate}
	\item Esse guia não é uma norma, é uma sugestão. Não tem como objetivo engessar a escrita. É uma sugestão de encadeamento dos itens inerentes ao processo de pesquisa, utilizando a formatação indicada pelo SiBI.
	\item Mais importante que o guia em si, é formar um pesquisador e deixar brotar seu estilo autoral a partir da convivência científica com seu orientador.  A elaboração do texto virá da interação orientador e orientando e deverá respeitar o estilo da comunidade acadêmica a que pertence.
	\item Esse documento foi elaborado dentre outras coisas a partir das diretivas do (SIBI UFRJ, 2025) e de (Silva, 2006) [Silva, 2025 - notas de aula da disciplina Metodologia de Pesquisa Científica ofertada em cursos lato e stricto sensu da UFRJ.]
	\item Esse guia não substitui a leitura e uso do Manual de formatação do SiBI.
	\item Esta versão do guia já está alinhada com a 9ª revisão do manual do SiBI \cite{ManualSibi2025}, de agosto de 2025.
	\item Caso esse material tenha sido útil tanto para a formatação quanto para o entendimento do que se espera de um trabalho acadêmico, e caso deseje, ficamos agradecidos com uma possível citação do nosso trabalho. Esse guia está contido dentro do que foi denominado Projeto Jacurutu \cite{Projeto_Jacurutu} e pode ser citado da mesma forma que a citação anterior e  cuja referência está presente na seção de Referências desse guia.
\end{enumerate}

Agora voltando ao objetivo desse guia, a seção introdutória versa sobre a apresentação inicial do trabalho acadêmico e deve indicar a motivação ou problemática, os objetivos da pesquisa, a relevância, a delimitação do assunto tratado e outros elementos necessários para situar o leitor.

Seja sucinto para não cansar o leitor, pois na introdução ele quer apenas uma ideia inicial do que virá à frente. Contudo, não deixe de “vender” a ideia, a relevância e as contribuições que emergirão do seu trabalho acadêmico nas próximas seções.

\section{Motivação ou problemática}

Relatar o que motivou a executar essa pesquisa, seja uma motivação da sua experiência pessoal ou uma motivação com base em problemas sociais atuais.

\section{Objetivos da pesquisa}

Nesta parte, você deve salientar qual o fenômeno a ser investigado e quais os objetivos principais e secundários ao pesquisar esse fenômeno.

Esta monografia visa responder qual pergunta? É importante que o aluno deixe claro a pergunta da pesquisa.

O objetivo pode ser organizado em Objetivo Geral e Objetivos Específicos, caso você queira falar do objetivo final e das partes que você precisará fazer para atingi-lo. Ou pode ser organizado em Objetivo Primário e Objetivos secundários, onde você fala do seu objetivo principal e de outros que você alcançará após a realização do primeiro.

\section{Relevância}

Por que é importante estudar esse tema? Qual a relevância acadêmica e prática de investigar tal assunto? 

A relevância fica diminuída quando apenas o próprio autor da monografia parece considerá-la importante. Assim, preferencialmente, a relevância deve ser levantada a partir da leitura de outros autores. Cite outros autores que consideram importante o estudo de tal tema.

\section{Delimitação da pesquisa}

A seção de delimitação de estudo visa demarcar o escopo de sua pesquisa, indicando ao leitor partes que não serão o alvo da sua investigação.

\section{Estrutura do trabalho}

O restante deste texto está organizado da seguinte forma. A Seção~\ref{cap:referencial_teorico} apresenta os conceitos básicos do referencial teórico necessário para um melhor entendimento do trabalho, com a apresentação dos assuntos tal, tal e tal. A Seção~\ref{cap:metodologia} relata qual foi a metodologia utilizada no trabalho. A Seção~\ref{cap:desenvolvimento} apresenta um detalhamento do  desenvolvimento do trabalho. Por fim a Seção~\ref{cap:conclusao} apresenta a conclusão, com um sumário dos resultados e sugestões para trabalhos futuros.